% GENERATED FILE. Edit canonical lessons, then run npm run generate:books.
% canonical-source-hash: fnv1a64:46f0093ebb7ad761
% canonical-lessons: GU-C126-pith, GU-C126-kamar, GU-C126-koni, GU-C126-hatheli, GU-C126-nakh

\chapter{Back, Waist, Elbow, Palm, Fingernail}
\label{ch:gu-back-waist-elbow-palm-fingernail}
\hlchaptermodality{126}

\begin{chapteropening}
\textbf{By the end of this chapter:} \emph{I can say a back, a waist, an elbow, a palm and a fingernail.}

Complete the last lesson of chapter 126: I can say a back, a waist, an elbow, a palm and a fingernail.
\end{chapteropening}

\section[\texorpdfstring{\gu{પીઠ} (pīṭh)}{પીઠ (pīṭh)}]{\gu{પીઠ} (pīṭh) — a back}

\label{lesson:GU-C126-pith}



\begin{quote}
Before the new one: say the Gujarati for a neck, then the Gujarati for a chest.
\end{quote}

\subsection*{You'll want to know: \gu{પીઠ}}
\textbf{\gu{પીઠ}} (\emph{pīṭh}) — \textquotedblleft{}a back\textquotedblright{}.

\begin{cousinweb}
\textbf{\gu{પીઠ}} (\emph{pīṭh}) is a tadbhava of Sanskrit \emph{pṛṣṭi}, \textquotedblleft{}vertebrae; back\textquotedblright{} (CDIAL 8370).
\end{cousinweb}

\begin{grammarlens}[title={What you've built}]
One more word for this chapter.
\end{grammarlens}

\subsection*{Your turn}
\begin{itemize}
\raggedright
  \item \emph{Say it:} \emph{pīṭh}
  \item \emph{Say it:} \emph{pīṭh}, once more
  \item \emph{Say it:} say \emph{pīṭh}
  \item \emph{Recall:} say the Gujarati for a neck, then the Gujarati for a chest, then say \emph{pīṭh} again
\end{itemize}

\begin{tcolorbox}[breakable,colback=teal!4,colframe=teal!35!black,title={Before you move on}]
What does \gu{પીઠ} mean? (\textquotedblleft{}A back\textquotedblright{}.) Say it once more.
\end{tcolorbox}
\section[\texorpdfstring{\gu{કમર} (kamar)}{કમર (kamar)}]{\gu{કમર} (kamar) — a waist}

\label{lesson:GU-C126-kamar}



\begin{quote}
Before the new one: say the Gujarati for a chest, then the Gujarati for a back.
\end{quote}

\subsection*{You'll want to know: \gu{કમર}}
\textbf{\gu{કમર}} (\emph{kamar}) — \textquotedblleft{}a waist\textquotedblright{}.

\begin{cousinweb}
\textbf{\gu{કમર}} (\emph{kamar}) is a loan, usually traced to Persian \emph{kamar}, \textquotedblleft{}waist\textquotedblright{}.
\end{cousinweb}

\begin{grammarlens}[title={What you've built}]
One more word for this chapter.
\end{grammarlens}

\subsection*{Your turn}
\begin{itemize}
\raggedright
  \item \emph{Say it:} \emph{kamar}
  \item \emph{Say it:} \emph{kamar}, once more
  \item \emph{Say it:} say \emph{kamar}
  \item \emph{Recall:} say the Gujarati for a chest, then the Gujarati for a back, then say \emph{kamar} again
\end{itemize}

\begin{tcolorbox}[breakable,colback=teal!4,colframe=teal!35!black,title={Before you move on}]
What does \gu{કમર} mean? (\textquotedblleft{}A waist\textquotedblright{}.) Say it once more.
\end{tcolorbox}
\section[\texorpdfstring{\gu{કોણી} (koṇī)}{કોણી (koṇī)}]{\gu{કોણી} (koṇī) — an elbow}

\label{lesson:GU-C126-koni}



\begin{quote}
Before the new one: say the Gujarati for a back, then the Gujarati for a waist.
\end{quote}

\subsection*{You'll want to know: \gu{કોણી}}
\textbf{\gu{કોણી}} (\emph{koṇī}) — \textquotedblleft{}an elbow\textquotedblright{}.

\begin{cousinweb}
\textbf{\gu{કોણી}} (\emph{koṇī}) is a tadbhava of the lexical Sanskrit \emph{kaphoṇi}, \textquotedblleft{}elbow\textquotedblright{} (CDIAL 2757). CDIAL reports that the Sanskrit word is more likely from Munda than from Dravidian.
\end{cousinweb}

\begin{grammarlens}[title={What you've built}]
One more word for this chapter.
\end{grammarlens}

\subsection*{Your turn}
\begin{itemize}
\raggedright
  \item \emph{Say it:} \emph{koṇī}
  \item \emph{Say it:} \emph{koṇī}, once more
  \item \emph{Say it:} say \emph{koṇī}
  \item \emph{Recall:} say the Gujarati for a back, then the Gujarati for a waist, then say \emph{koṇī} again
\end{itemize}

\begin{tcolorbox}[breakable,colback=teal!4,colframe=teal!35!black,title={Before you move on}]
What does \gu{કોણી} mean? (\textquotedblleft{}An elbow\textquotedblright{}.) Say it once more.
\end{tcolorbox}
\section[\texorpdfstring{\gu{હથેળી} (hatheḷī)}{હથેળી (hatheḷī)}]{\gu{હથેળી} (hatheḷī) — a palm}

\label{lesson:GU-C126-hatheli}



\begin{quote}
Before the new one: say the Gujarati for a waist, then the Gujarati for an elbow.

Ask \textquotedblleft{}why?\textquotedblright{}, then begin the answer. (\emph{Kem?} … \emph{kemke} …)
\end{quote}

\subsection*{You'll want to know: \gu{હથેળી}}
\textbf{\gu{હથેળી}} (\emph{hatheḷī}) — \textquotedblleft{}a palm\textquotedblright{}.

\begin{cousinweb}
\textbf{\gu{હથેળી}} (\emph{hatheḷī}) goes back to Sanskrit \emph{hastatala}, \textquotedblleft{}palm of the hand\textquotedblright{}: \emph{hasta}, \textquotedblleft{}hand\textquotedblright{}, and \emph{tala}, \textquotedblleft{}surface\textquotedblright{} (CDIAL 14029). CDIAL lists Hindi \emph{hathelī} there, though no Gujarati form.
\end{cousinweb}

\begin{grammarlens}[title={What you've built}]
One more word for this chapter.
\end{grammarlens}

\subsection*{Your turn}
\begin{itemize}
\raggedright
  \item \emph{Say it:} \emph{hatheḷī}
  \item \emph{Say it:} \emph{hatheḷī}, once more
  \item \emph{Say it:} say \emph{hatheḷī}
  \item \emph{Recall:} say the Gujarati for a waist, then the Gujarati for an elbow, then say \emph{hatheḷī} again
\end{itemize}

\begin{tcolorbox}[breakable,colback=teal!4,colframe=teal!35!black,title={Before you move on}]
What does \gu{હથેળી} mean? (\textquotedblleft{}A palm\textquotedblright{}.) Say it once more.
\end{tcolorbox}
\section[\texorpdfstring{\gu{નખ} (nakh)}{નખ (nakh)}]{\gu{નખ} (nakh) — a fingernail}

\label{lesson:GU-C126-nakh}



\begin{quote}
Before the new one: say the Gujarati for an elbow, then the Gujarati for a palm.
\end{quote}

\subsection*{You'll want to know: \gu{નખ}}
\textbf{\gu{નખ}} (\emph{nakh}) — \textquotedblleft{}a fingernail\textquotedblright{}.

\begin{cousinweb}
\textbf{\gu{નખ}} (\emph{nakh}) is Sanskrit \emph{nakha}, \textquotedblleft{}nail of finger or toe, claw\textquotedblright{}, taken over whole as a learned loan, a tatsama (CDIAL 6914). CDIAL lists the inherited Old Gujarati form as \emph{naha}.
\end{cousinweb}

\begin{grammarlens}[title={What you've built}]
One more word for this chapter.
\end{grammarlens}

\subsection*{Your turn}
\begin{itemize}
\raggedright
  \item \emph{Say it:} \emph{nakh}
  \item \emph{Say it:} \emph{nakh}, once more
  \item \emph{Say it:} say \emph{nakh}
  \item \emph{Recall:} say the Gujarati for an elbow, then the Gujarati for a palm, then say \emph{nakh} again
\end{itemize}

\begin{tcolorbox}[breakable,colback=teal!4,colframe=teal!35!black,title={Before you move on}]
What does \gu{નખ} mean? (\textquotedblleft{}A fingernail\textquotedblright{}.) Say it once more.
\end{tcolorbox}
